\section{アルゴリズムの説明}

本節では、各課題を解くために採用したアルゴリズムを、プログラミング言語に依存しない形で説明する。

\subsection{動画の読み込み(課題 1)}

動画ファイルは、連続する静止画(フレーム)の列と、その再生速度などのメタデータから
構成される。動画コンテナには幅・高さ・1 秒あたりのフレーム数(FPS)・総フレーム数が
メタデータとして格納されているため、動画をデコーダで開き、これらのメタデータを
問い合わせて出力すればよい。

\subsection{サムネイル画像の作成(課題 2)}

サムネイルは動画の内容を 1 枚で代表する画像である。先頭フレームはフェードインや
無内容な場面であることが多く、代表性が保証されない。そこで本課題では
「動画全体の明るさ分布に最も近い明るさ分布を持つフレームが、動画を最もよく代表する」
という基準を採用した。手順は以下のとおりである。

\begin{enumerate}
\item 動画から一定間隔でフレームをサンプリングする。
\item 各フレームをグレースケール化し、画素値の正規化ヒストグラム
      $\boldsymbol{h}_i$ を計算する。
\item 全サンプルの平均ヒストグラム $\bar{\boldsymbol{h}}$ を求める。
\item 各フレームについて平均ヒストグラムとのユークリッド距離を計算する。
\begin{equation}\label{eq:thumb}
d_i = \| \boldsymbol{h}_i - \bar{\boldsymbol{h}} \|_2
\end{equation}
\item $d_i$ が最小のフレームを代表フレームとして選択する。
\end{enumerate}

式 \ref{eq:thumb} の $d_i$ が小さいフレームは動画の平均的な画面構成に近く、
場面転換の途中や極端に明るい・暗い例外的なフレームが選ばれることを避けられる。

\subsection{ファイル名のパース(課題 3)}

ファイル名は「タイトル - ID.拡張子」という規則で構成されている。したがって、
(1) まず拡張子を取り除き、(2) 残った文字列を区切り文字列「{\tt "\ -\ "}」
(空白・ハイフン・空白)で分割すれば、前半がタイトル、後半が ID となる。
タイトル自体にハイフンが含まれる場合を考慮し、分割は右端の区切りで 1 回だけ行う。
これにより「A-B - 12345.mp4」のようなファイル名でもタイトル「A-B」と ID「12345」に
正しく分けられる。

\subsection{色反転と動画の保存(課題 4)}

画素値が 8 ビット(最大値 255)で表されるとき、色反転は各画素・各チャンネルの値
$p$ を次式で置き換える操作である。

\begin{equation}\label{eq:invert}
p' = 255 - p
\end{equation}

式 \ref{eq:invert} により、暗い画素($p$ が小さい)は明るく($p'$ が大きく)、
明るい画素は暗く変換される。動画に対しては、全フレームに式 \ref{eq:invert} を適用し、
縦・横を 0.3 倍に縮小したうえで、元動画と同じ FPS で動画ファイルとして書き出す。

\subsection{音声の可視化(課題 5)}

音声信号は時間領域の 1 次元信号であり、そのままでは周波数方向に埋め込まれた
メッセージを読み取れない。そこで短時間フーリエ変換
(STFT: Short-Time Fourier Transform)を用いる。STFT は信号を長さ $N$ の窓で
切り出しながら、窓ごとに離散フーリエ変換を行う操作であり、信号 $x[n]$、窓関数 $w[n]$
に対して次式で定義される。

\begin{equation}\label{eq:stft}
X(m, k) = \sum_{n=0}^{N-1} x[n + mH]\, w[n]\, e^{-j 2 \pi k n / N}
\end{equation}

ここで $m$ は窓の位置(時間方向)、$k$ は周波数ビン、$H$ はホップ幅
(窓をずらす幅)である。$|X(m,k)|^2$ を時間-周波数平面に濃淡で描いたものが
スペクトログラムであり、周波数方向に描かれた文字列を画像として可視化できる。
窓長 $N$ を大きくすると周波数分解能が上がる代わりに時間分解能が下がる
(不確定性関係)ため、メッセージが読みやすくなるよう $N$ とホップ幅を調節する。
また、パワーは対数(dB)スケール
\begin{equation}\label{eq:db}
P_{\mathrm{dB}} = 10 \log_{10} |X(m,k)|^2
\end{equation}
で表示することで、振幅差の大きい成分が共存していても濃淡の判別ができるようにする。

\subsection{n-gram(課題 6)}

\subsubsection{2-gram の抽出と頻度集計(6-a)}

2-gram とは、文章を単語単位で区切ったときの連続する 2 語の組である。
単語列 $w_1, w_2, \ldots, w_M$ からは $M-1$ 個の 2-gram
$(w_1,w_2), (w_2,w_3), \ldots, (w_{M-1},w_M)$ が得られる。
テキストを (1) 小文字化し、(2) アルファベット列のみを単語として切り出し、
(3) 隣接する 2 語の組を列挙し、(4) 組ごとの出現回数を数え、
(5) 出現回数の降順に並べて上位 10 個を棒グラフとして描画する。

\subsubsection{stop words 除去(6-b)}

stop words は冠詞・be 動詞・前置詞・接続詞・代名詞など、出現頻度は高いが
文章の内容的特徴を担わない語である。単語列から stop words を取り除いた後に
2-gram を構成することで、内容語同士の連接だけが集計対象となり、
文章の主題を反映した 2-gram が上位に現れることが期待される。
除去は「単語列をフィルタしてから 2-gram を作る」方式とした。このため、
元の文章では stop words を挟んで離れていた 2 つの内容語が新たに隣接する
2-gram として数えられることがある。この設計の影響は考察で述べる。

\subsection{動画の変わり目検出(課題 7)}

動画の変わり目(ショット境界)では、隣接フレーム間で画面全体の画素値が大きく変化する。
そこで、フレーム $t$ のグレースケール画像を $I_t(x, y)$ とし、隣接フレーム間の
変化量 $D(t)$ を次式で定義する。

\begin{equation}\label{eq:sad}
D(t) = \sum_{x} \sum_{y} \left| I_t(x, y) - I_{t-1}(x, y) \right|
\end{equation}

同一ショット内ではフレーム間の変化は被写体やカメラの動きによる小さな値に留まるが、
ショット境界では画面全体が入れ替わるため $D(t)$ が突出して大きくなる。
したがって $D(t)$ がある閾値 $\theta$ を超えるフレームを変わり目と判定できる。
閾値の決め方として次の 3 方式を用いた。

\begin{enumerate}
\item \textbf{統計量に基づく方式}: $D(t)$ の平均 $\mu$ と標準偏差 $\sigma$ から
\begin{equation}\label{eq:thresh}
\theta = \mu + k \sigma
\end{equation}
とする。変わり目は全フレームのごく一部なので、$D(t)$ の分布はショット内変化を
中心に集中し、変わり目は外れ値として $\mu + k\sigma$ の外側に現れる。
\item \textbf{上位 $n$ ピーク方式}: 変わり目の個数(本課題では 3)が既知であることを
利用し、$D(t)$ の大きい順に $n$ 個のフレームを選ぶ。
\item \textbf{移動平均基準方式}: $D(t)$ の近傍 $W$ フレームの移動平均
$\bar{D}(t)$ を局所的な基準線とし、$D(t) > k' \bar{D}(t)$ となるフレームを選ぶ。
動きの多い区間と静かな区間が混在しても、局所的な変化の水準に応じて閾値が
自動的に追従する。
\end{enumerate}

\subsection{オブジェクト指向設計(課題 8)}

課題 2・3・4 の処理は、いずれも「1 本の動画」に付随する操作である。そこで動画 1 本を
表す Video クラスを定義し、動画のパス・タイトル・ID・サムネイルをインスタンス変数として
保持し、各処理をメソッドとして所属させる。これにより、データ(動画の属性)と
それを操作する処理が 1 つの単位にまとまり、複数の動画を扱う場合も動画ごとに
インスタンスを作るだけでよくなる。

さらに、動画処理を「フレーム 1 枚を変換する処理」と「全フレームへの適用と入出力」に
分離する。フレーム変換関数(色反転・グレースケール化・二値化)を引数として受け取る
高階関数方式にすることで、新しい画像処理を追加するときも動画入出力のコードを
複製せず、フレーム変換関数を 1 つ書いて渡すだけで済む。
